\section{极限问题的提出}\label{sec:3.1}

17 世纪建立的微积分，已有了较为一般的概念，更重要的是提供了一套关于连续变量的崭新的算法，从而成为解决天文、力学、光学与技术问题的重要工具．但那时的微积分，只能说是直观的微积分．因为它的一些基本概念并不确切，甚至逻辑推导前后矛盾．例如，微积分的一个最重要的概念——微商，它反映了一个连续量对另一个连续量的``瞬时变化率''，其概念就很不确切，逻辑推导前后矛盾．牛顿当时把连续量称为流量，把变化率称为流数（把自变量理解为时间）．例如把自由落体的路程看作流量，流数就是瞬时速度．用今天的符号，牛顿当时推导流数的方法是：对自由落体的运动规律 $s = \frac{1}{2}gt^2$，给时间 $t$ 一个微小的改变 $h$（他称之为瞬——moment），则平均速度为
\[
\frac{\frac{1}{2}g(t+h)^2-\frac{1}{2}gt^2}{h} = gt + \frac{1}{2}gh.
\]
然后令 $h=0$，便得瞬时速度 $gt$．凭着类似的推导，微积分的演算系统建立起来了．但这里的瞬 $h$ 究竟是什么？贝克莱（Berkeley，1685---1753）就曾指出这里有逻辑错误．因为在计算平均速度时，要用 $h$ 去除，必须假设 $h \neq 0$．但为最后得到瞬时速度，必须令 $h=0$，才能算出 $gt$．$h$ 非零又是零，$h$ 到底是什么？难怪贝克莱讥讽说这是``消失的量的鬼魂''．这个问题困惑了数学家很长时间．达朗贝尔认为说清楚这点，应该用极限，但他并没有把它说清楚．拉格朗日企图用代数推演的方法逃避这个困难，但这脱离了微积分的物理意义，很难应用．直到 19 世纪柯西才真正用极限的概念把它基本说清楚，而魏尔斯特拉斯最终用 $\varepsilon - \delta$ 的语言，彻底解决了这个困难，从而推动了近代分析的蓬勃发展．

极限最早的观念，在国外有所谓``穷竭法''，在中国有所谓割圆术，即把圆近似地割成边数很多、边长很细的正多边形来计算圆面积．魏晋时代的刘徽就说过：``割之弥细，所失弥小，割之又割，以至于不可割，则与圆周合体而无所失矣''．后来，和微积分同时产生而用得很多的一种运算，就是无穷多项相加
\[
a_1 + a_2 + \cdots + a_n + \cdots.
\]
怎样理解这无穷多项相加？这只能理解为有限项相加，但越加越多，最后看它的变化趋势，这就是所谓的数列极限．把数列极限的思想用来观察求瞬时速度的推导，就是把 $h$ 看成越变越小而且无限地变小（但不为零），来研究平均速度的变化趋势，从而建立一套从平均速度求瞬时速度的计算方法，这就是所谓的函数极限．今天，我们学习微积分，当然不必重复历史的过程，而可以从严格的极限概念开始，严格地讲述微积分的概念、演算及应用．

下面我们先从数列极限入手，然后讲述函数极限，再用函数极限，给出函数连续的概念，并最后证明，一切初等函数在其定义域是连续的．从下一章开始转入微积分．
